\begin{textsample}{NEG Dim 7 – global\_south\_2025\_06 – Score -5.00 – global\_south\_2025\_06/126471584.txt}  \label{ex:f7_neg_280}
British punk-rap duo Bob Vylan is facing police scrutiny after leading Glastonbury \textbf{Festival} \textbf{crowds} in chants of"Free, free Palestine"and"Death, death to the IDF"during a Saturday performance.
Avon and Somerset Police confirmed they are reviewing footage to assess whether any criminal offences were committed that warrant an investigation.
British police officers are also examining comments by the Irish rap trio Kneecap, whose members have also been highly critical of Israel and its military assault on Gaza.
The incident quickly drew criticism from officials and pro-Israel groups.
The Israeli Embassy in the UK called the chants"inflammatory and hateful,"while the UK Health Secretary Wes Streeting labelled the performance"appalling"in an interview with Sky News.
He also questioned how such content was broadcast live by the BBC to millions of viewers.
UK Prime Minister Keir Starmer on Sunday added his voice to those condemning the punk-rap group.
He told The Telegraph Sunday that"there is no excuse for this kind of appalling @ @ @ @ @ @ @ @ @ @ should not be given a platform, and that goes for any other performers making threats or inciting violence,"he added."The BBC needs to explain how these scenes came to be broadcast,"he said, referring to the country ' s national broadcaster.
One of Kneecap ' s members wore a T-shirt dedicated to the Palestine Action Group, which is about to be banned under UK terror laws.
The \textbf{festival} ' s organisers said Bob Vylan ' s comments had"very much crossed a line very much"."We are urgently reminding everyone involved in the production of the \textbf{festival} that there is no place at Glastonbury for antisemitism, \textbf{hate} speech or incitement to violence,"the \textbf{festival} said in a statement.
Avon and Somerset police said Saturday that video evidence would be assessed by officers"to determine whether any offences may have been committed that would require a criminal investigation".
The chants about Israel ' s military were led by Bob Vylan ' s frontman Bobby Vylan, and were broadcast live @ @ @ @ @ @ @ @ @ @ most popular \textbf{music} \textbf{festival}."I thought it ' s appalling,"Streeting said of the chants."I think the BBC and Glastonbury have got questions to answer about how we saw such a spectacle on our screens,"he told Sky News.
The Israel embassy said in a statement late Saturday that it was"deeply disturbed by the inflammatory and hateful rhetoric expressed on stage at the Glastonbury \textbf{Festival}".
But Streeting also aimed the embassy, telling it to"get your own house in order"."I think there ' s a serious point there by the Israeli embassy.
I wish they ' d take the violence of their own citizens towards Palestinians more seriously,"he said, citing Israeli settler violence in the West Bank.
Festival-goer Joe McCabe, 31, told AFP that while he did not necessarily agree with Vylan ' s statement,"I certainly think the message of questioning what ' s going on there ( in Gaza ) is right."A spokesperson for the @ @ @ @ @ @ @ @ @ @ and the broadcaster had"no plans"to make the performance available on its on-demand service.
While the BBC reported that Kneecap ' s set was not broadcast live due to editorial concerns surrounding impartiality, it announced on Sunday that an edited version was made available on iPlayer.
It said the edits ensured the content fell within the"limits of artistic expression"in line with its editorial guidelines, and they had put warnings for strong language in the video.
Kneecap, which has made headlines in recent months with its pro-Palestinian and anti-Israel stance, also led \textbf{crowds} in chanting abuse against UK Prime Minister Starmer.
Starmer and other politicians had said the band should not perform after its member Liam O'Hanna, known by his stage name Mo Chara, was charged with a terror offence.
He appeared in court this month, accused of having displayed a Hezbollah flag while saying"Up Hamas, Up Hezbollah"after a video resurfaced of a London \textbf{concert} last year.
Both Hezbollah and Hamas @ @ @ @ @ @ @ @ @ @ offence to express support for them.
O'Hanna has denied the charge and told the Guardian newspaper in an interview published Friday that"it was a joke -- we ' re playing characters".
The group, which regularly leads \textbf{crowds} in chants of"Free Palestine"during its \textbf{concerts}, apologised this year after a 2023 video emerged appearing to show one singer calling for the death of British Conservative lawmakers.
Israel ' s military assault on Gaza has killed at least 56,412 people, mostly civilians, according to Gaza ' s health ministry.
The United Nations considers these figures to be reliable.
Comments 1000 Characters MahmoodJun 30, 2025 01:14pm Please someone remind Kier Stammer of the Britain ' s role in the creation of the fascist apartheid State of Israel with Balfour Declaration, giving the Jews the right to Palestinian land.
The crimes by committed by the Zionists in Palestine since 1948 are the direct results of British and German complicity to transfer the problem to someone else.

% matched lemmas: concert, crowd, festival, hate, music
\end{textsample}
